% 妊神星（综述）
% license CCBYSA3
% type Wiki

本文根据 CC-BY-SA 协议转载翻译自维基百科\href{https://en.wikipedia.org/wiki/Haumea}{相关文章}。

Haumea（小行星编号：136108 Haumea）是一颗位于海王星轨道之外的矮行星。$^{[25]}$
它由加州理工学院（Caltech）的 Mike Brown 领导的团队于 2004 年在帕洛马天文台发现；随后，西班牙内华达山脉天文台 José Luis Ortiz Moreno 领导的团队在 2005 年正式宣布了这一发现，他们在 2003 年拍摄的先前图像中识别出了该天体。从那次宣布起，它获得了临时编号 2003 EL61。

2008 年 9 月 17 日，它被正式命名为 Haumea，源自夏威夷生育与分娩女神。国际天文学联合会（IAU）基于其可能具备矮行星地位的预期而赋予这一名称。根据标称估计，它是仅次于厄里斯和冥王星的第三大已知海王星外天体，其大小约等同于天王星的卫星泰坦妮娅。Haumea 的最早“预发现”（precovery）影像可以追溯至 1955 年 3 月 22 日。$^{[9]}$

Haumea 的质量约为冥王星的三分之一、地球的 1/1400。虽然尚未直接观测到其形状，但根据其光变曲线的计算结果，与 雅可比椭球体（Jacobi ellipsoid）一致（这是一种在满足矮行星条件时会呈现的形状），其长轴大约是短轴的两倍。

2017 年 10 月，天文学家宣布在 Haumea 周围发现了一套环系统——这是首个在海王星外天体以及矮行星上发现的环系统。

Haumea 的引力直到最近仍被认为足以使其达到流体静力平衡，但目前这一点仍不确定。Haumea 修长的形状、快速自转、环结构以及高反照率（源自其晶质水冰表面），被认为都是一次巨大碰撞的结果。这场碰撞使 Haumea 成为一个碰撞族群（Haumea 家族）的最大成员，该族群包括多个大型海王星外天体，以及 Haumea 的两个已知卫星：Hiʻiaka 与 Namaka。
\subsection{历史}
\subsubsection{发现}
由加州理工学院（Caltech）的 Mike Brown、耶鲁大学的 David Rabinowitz，以及位于夏威夷双子天文台的 Chad Trujillo 组成的团队，于 2004 年 12 月 28 日在他们拍摄的2004 年 5 月 6 日的图像中发现了 Haumea。2005 年 7 月 20 日，他们在网上发布了一篇摘要，准备在 2005 年 9 月的一次会议上正式宣布这一发现。$^{[26]}$

在此期间，西班牙内华达山脉天文台安达卢西亚天体物理研究所（Instituto de Astrofísica de Andalucía）的José Luis Ortiz Moreno及其团队，在2003 年 3 月 7–10 日拍摄的图像中找到了 Haumea。2005 年 7 月 27 日晚，Ortiz 向小行星中心（MPC）发送了关于他们发现的邮件。$^{[27]}$

Brown 最初承认发现权应归 Ortiz 团队，$^{[28]}$ 但在得知西班牙团队在官方发现宣布的前一天访问过 Brown 的观测日志后，开始怀疑对方存在欺诈行为，而 Ortiz 团队在宣布中并未像惯例那样披露这一事实。这些日志包含足够的数据，使 Ortiz 团队能够在其 2003 年的图像中提前定位（precover）Haumea，并且在 Ortiz 在7 月 29 日安排望远镜时间以获取确认图像、并向 MPC 再次提交发现报告之前，又再次访问了这些日志。Ortiz 后来承认访问过 Caltech 的观测日志，但否认任何不当行为，声称自己只是核实他们是否发现了一个新的天体。$^{[29]}$

根据 IAU（国际天文学联合会）规定，小行星的发现权归属于最先向 MPC（小行星中心）提交足够轨道数据以确定其轨道的团队，并且获得发现权的团队有权优先命名该天体。
然而，IAU 在2008 年 9 月 17 日发布的关于 Haumea 命名的官方公告（该命名由专门负责预期为矮行星的天体设立的联合委员会决定）中，并未提及发现者。公告中将发现地点列为西班牙团队的内华达山脉天文台，$^{[30][31]}$ 但最终采用的名称“Haumea”却是 Caltech 团队的提议。

Ortiz 团队曾提议命名为“Ataecina”，取自古伊比利亚春之女神；$^{[27]}$ 但由于该女神属于冥界神（chthonic deity），这一名字更适合用于命名“冥卫星类（plutino）”天体，而 Haumea 并非此类天体，因此未被采用。
\subsubsection{名称与符号}
在被正式命名之前，加州理工学院（Caltech）发现团队内部给 Haumea 起了绰号“Santa（圣诞老人）”，因为他们是在2004年12月28日，也就是圣诞节刚过后发现该天体的。$^{[32]}$

西班牙团队于2005年7月首次向小行星中心（MPC）提交了发现申请。2005 年 7 月 29 日**，Haumea 根据西班牙团队发现图像的日期，获得了临时编号2003 EL61。2006 年 9 月 7 日**，它被编号并正式列入小行星目录，成为(136108) 2003 EL61。

根据当时国际天文学联合会（IAU）规定：经典柯伊伯带天体应以与创造相关的神话人物命名，$^{[33]}$ 因此 2006 年 9 月，Caltech 团队向 IAU 正式提交了来自夏威夷神话的名称建议，用于命名 (136108) 2003 EL61 及其卫星，以“向卫星发现地致敬”。$^{[34]}$ 这些名称由 Caltech 团队成员 David Rabinowitz 提出。$^{[25]}$

Haumea 是夏威夷岛（Hawaiʻi）的守护女神，双子座天文台（Gemini）与凯克天文台（W. M. Keck Observatory）就坐落在冒纳凯亚山（Mauna Kea）。此外，她也被视作大地女神 Papa，是天空之神 Wākea（象征太空）的妻子。$^{[35]}$ 在当时，这一命名被认为非常贴切，因为当时 Haumea 被认为几乎完全由岩石构成，而不像其他柯伊伯带天体那样拥有“岩石核心 + 厚冰幔”的结构。$^{[36][37]}$

此外，Haumea 是 生育与分娩女神，据说许多子女从她身体的不同部位诞生；$^{[35]}$ 这与理论推测相对应——Haumea 主体在远古一次碰撞中碎裂，形成了大量冰质天体（即所谓 Haumea 碰撞家族）。$^{[37]}$ Haumea 已知的两颗卫星也被认为是在这一碎裂事件中形成，因此以其两个女儿 Hiʻiaka 与 Nāmaka 命名。$^{[36]}$

西班牙 Ortiz 团队提出的命名 Ataecina 不符合 IAU 命名规定，因为冥界神（chthonic deity）的名字仅用于稳定共振的柯伊伯带天体，例如与海王星呈 3:2 共振的冥卫星类（plutinos），而 Haumea 处于不稳定的 7:12 间歇共振状态，因此严格意义上并非共振天体。此命名规则在 2019 年末 被 IAU 明确：冥界神名称今后专属用于冥卫星类天体。

Haumea 的行星符号为 ⟨🝻⟩，现收录于 Unicode 编码 U+1F77B。$^{[38]}$行星符号在现代天文学中已较少使用，🝻 主要由占星师使用，$^{[39]}$ 但 NASA 也曾使用过该符号。$^{[40]}$
该符号由马萨诸塞州软件工程师 Denis Moskowitz 设计，其造型结合并简化了夏威夷岩画中代表“女性”和“分娩”的符号。$^{[41]}$
\subsection{轨道}
Haumea 的轨道周期为 284 个地球年，近日点距离约 35 天文单位（AU），轨道倾角为 28°。$^{[9]}$ 它在 1992 年初经过远日点，目前距离太阳已经超过 50 AU。$^{[23]}$ Haumea 将在 2133 年抵达近日点。$^{[10]}$

Haumea 的轨道偏心率略高于其碰撞家族其他成员。人们认为，这主要是因为 Haumea 与海王星之间存在一个较弱的 7:12 轨道共振，在大约十亿年的时间尺度上逐渐改变了其原始轨道，$^{[37][42]}$ 这种演化过程受 Kozai 效应（古在效应）影响，该效应允许轨道倾角与偏心率之间发生交换。$^{[37][43][44]}$

Haumea 的视星等为 17.3，$^{[23]}$ 是柯伊伯带中继冥王星与鸟神星（Makemake）之后的第三亮天体，使用大型业余望远镜即可轻松观测。$^{[45]}$

然而，由于行星与大多数小型太阳系天体都形成于原始太阳星云盘，因此它们具有共同的轨道平面投影，即黄道面（ecliptic）。$^{[46]}$ 早期对远距天体的搜索，大多集中在天空上接近黄道面的区域。当这一地区逐渐被充分勘测后，后续的天空巡天开始寻找轨道倾角更高、以及更遥远、运动更缓慢的天体。$^{[47][48]}$ 正是在这些扩展后的巡天中，才覆盖到 Haumea 所在的天空位置，因此得以发现该天体，因为它具有较高的轨道倾角，并且目前位于远离黄道的位置。
\subsubsection{可能的与海王星轨道共振}
哈乌梅亚被认为处于与海王星的间歇性 7:12 轨道共振$^{[37]}$。其升交点 $\Omega$ 的进动周期约为 460 万年，每一个进动周期共振会中断两次（即每约 230 万年一次），在约十万年后又重新恢复$^{[5]}$。由于这种共振并不稳定，Marc Buie 将其归类为非共振状态$^{[49]}$。
\subsection{自转}
哈乌梅亚在约 3.9 小时的周期内呈现出剧烈的亮度变化，这只能用其自转周期来解释$^{[50]}$。这比太阳系中任何其他已知达到流体静力平衡（equilibrium body）的天体都更快，事实上也比任何直径超过 100 公里的已知天体自转得更快$^{[45]}$。大多数达到平衡形状的自转天体会被压扁成扁球体（oblate spheroid），但哈乌梅亚自转速度极快，因此被拉伸成一个三轴椭球体。如果哈乌梅亚再快一些，它将变形为“哑铃形”，最终可能裂成两个天体$^{[25]}$。这种快速自转被认为是形成其卫星及碰撞族群的大撞击事件所造成的$^{[37]}$。

目前从地球观察，哈乌梅亚的赤道平面几乎是边缘视角（edge-on），并且相对于其环轨道和最外侧卫星希伊阿卡（Hiʻiaka）的轨道平面略有偏移。虽然 Ragozzine 和 Brown 在 2009 年最初假设二者是共面，但他们关于哈乌梅亚卫星碰撞形成模型的研究一直表明，哈乌梅亚赤道平面至少与希伊阿卡的轨道平面对齐，差异约为 1°$^{[15]}$。

